%% kapa: exam=1 round=37 year=2026 period=1 form=A subject=civil_law subject_ko=민법 seq=24 number=24
%% source: Q-Net 감정평가사 기출문제 (cst00309) artlSeq=5258924 file=raw/1008/5258924/2026년 제37회 감정평가사 1차 1교시.pdf
%% conversion: kordoc 4.24.51 -> tools/gen1.py (auto); review=proofread-sonnet; flags=none
\begin{question}{24}
등기를 마치지 않더라도 물권변동의 효력이 발생하는 경우가 아닌 것은? (다툼이 있으면 판례에 따름)
\begin{choices}
  \item 법정대위에 의해 저당권이 이전하는 경우
  \item 공유자 사이에 공유토지분할에 관한 조정이 성립한 경우
  \item 법정저당권을 취득한 경우
  \item 재결수용에 의해 부동산 소유권을 취득한 경우
  \item 부동산의 공경매에서 매수인이 매각대금을 완납한 경우
\end{choices}
\end{question}
\answer{2}
